% sections/02b_theory_inverse.tex  (TV7)

\subsection{Determinant, inverse matrix and invertible filters}\label{sec:theory-inverse}\writer{TV7}

As explained in Section~\ref{sec:theory-filters}, a color filter maps $p$ to
$p' = Mp$. In Task~12 the project asks if we can get the original colors back
after the filter is applied. To answer this, we need the inverse matrix and the
determinant.

\paragraph{Inverse matrix.}
A square matrix $A$ is \emph{invertible} if there is a matrix $A^{-1}$ such that
$AA^{-1} = A^{-1}A = I$, where $I$ is the identity matrix. If there is no such
matrix, $A$ is \emph{singular}. If the filter matrix $M$ is invertible, we can
multiply $p' = Mp$ by $M^{-1}$ and get $M^{-1}p' = p$. So $M^{-1}$ is the inverse
filter: it brings every color back.

\paragraph{When is a matrix invertible?}
For an $n \times n$ matrix $A$, the following statements are equivalent
\cite{strang2016}:
\begin{itemize}
  \item $A$ is invertible;
  \item $\det(A) \neq 0$;
  \item the rows (or the columns) of $A$ are linearly independent;
  \item $Ax = 0$ has only the solution $x = 0$.
\end{itemize}
For a filter, the last statement means that two different colors never give the
same result. If they did, we could not know which one was the original.

\paragraph{How to compute.}
For a $3 \times 3$ matrix, the determinant can be found by expanding along the
first row:
\[
  \det\begin{pmatrix} a & b & c \\ d & e & f \\ g & h & i \end{pmatrix}
  = a(ei - fh) - b(di - fg) + c(dh - eg)
\]
The inverse is $A^{-1} = \frac{1}{\det(A)}\,C^{T}$, where $C$ is the matrix of
cofactors $C_{ij} = (-1)^{i+j} D_{ij}$ and $D_{ij}$ is the determinant of the
$2 \times 2$ matrix left after deleting row $i$ and column $j$. We must divide by
$\det(A)$, so the inverse does not exist when $\det(A) = 0$. In our code we use
\texttt{np.linalg.inv}, which uses Gaussian elimination instead of this formula
because it is faster and more accurate \cite{higham2002}.

\paragraph{Example by hand: \texttt{UserMatrix}.}
The filter in Task~11 is
\[
  U = \begin{pmatrix}
    0.7  & 0.15 & 0.15 \\
    0.15 & 0.7  & 0.15 \\
    0.15 & 0.15 & 0.7
  \end{pmatrix}
\]
The cofactors of the first row are
\[
  C_{11} = 0.7 \cdot 0.7 - 0.15 \cdot 0.15 = 0.4675, \qquad
  C_{12} = C_{13} = -(0.15 \cdot 0.7 - 0.15 \cdot 0.15) = -0.0825
\]
so
\[
  \det(U) = 0.7 \cdot 0.4675 + 0.15 \cdot (-0.0825) + 0.15 \cdot (-0.0825)
          = 0.32725 - 0.02475 = 0.3025 \neq 0
\]
The determinant is not zero, so $U$ is invertible. This is the same number that
our code prints in Task~12.

Because all diagonal entries of $U$ are equal and all off-diagonal entries are
equal, all cofactors on the diagonal are $0.4675$ and all the others are
$-0.0825$. Dividing them by $\det(U)$ gives
\[
  \frac{0.4675}{0.3025} = \frac{17}{11} \approx 1.5455, \qquad
  \frac{-0.0825}{0.3025} = -\frac{3}{11} \approx -0.2727
\]
Therefore
\[
  U^{-1} = \frac{1}{11}
  \begin{pmatrix} 17 & -3 & -3 \\ -3 & 17 & -3 \\ -3 & -3 & 17 \end{pmatrix}
  \approx
  \begin{pmatrix}
     1.5455 & -0.2727 & -0.2727 \\
    -0.2727 &  1.5455 & -0.2727 \\
    -0.2727 & -0.2727 &  1.5455
  \end{pmatrix}
\]
To check, the first row of $U$ times the first column of $U^{-1}$ is
$(0.7 \cdot 17 - 0.15 \cdot 3 - 0.15 \cdot 3)/11 = 11/11 = 1$.
This matches \texttt{InvUserMatrix} in our code.

\paragraph{What the inverse filter does (Q5).}
In each row of $U$ the numbers add up to 1, so a gray pixel $(g,g,g)$ does not
change. Other colors get $15\%$ from the two other channels, so they move closer
to gray and the image looks faded. The rows of $U^{-1}$ also add up to 1, but the
numbers outside the diagonal are negative. So $U^{-1}$ takes away part of the
other channels and moves colors away from gray. When we apply it to
\texttt{ImJPG\_User}, we get the original image back. The difference from
\texttt{ImJPG} is at most 1, because \texttt{ImJPG\_User} was rounded to whole
numbers. When we apply $U^{-1}$ to the original \texttt{ImJPG}, the colors become
stronger. Some values go below 0 or above 255 and are clipped, so this result
cannot be undone exactly. Figure~\ref{fig:tv7-task12} shows these results.

\begin{figure}[htbp]
  \centering
  \includegraphics[width=0.95\linewidth]{figures/TV7_task12.png}
  \caption{Task 12: the original image, the image after \texttt{UserMatrix}
    (faded colors), the image recovered with \texttt{InvUserMatrix}, and the
    inverse filter applied to the original image (stronger colors).}
  \label{fig:tv7-task12}
\end{figure}

\paragraph{Singular filters lose information.}
If $\det(M) = 0$, there is a vector $x \neq 0$ with $Mx = 0$, so the colors $p$
and $p + x$ give the same result. Three filters in this project are singular:
\begin{itemize}
  \item \texttt{GrayMatrix} (all entries are $1/3$) keeps only the average
    $(R+G+B)/3$. For example, the two different colors $(255,0,0)$ (red) and
    $(0,255,0)$ (green) both become $(85,85,85)$. Looking at the gray pixel, we
    cannot know if the original color was red or green.
  \item \texttt{RedMatrix} keeps only the red channel, so green and blue are
    lost. For example, $(100,50,200)$ and $(100,0,0)$ both become $(100,0,0)$.
  \item \texttt{DeleteGreenMatrix} (Task~8) keeps red and blue but deletes green.
    The colors $(200,30,50)$ and $(200,220,50)$ both become $(200,0,50)$, see
    Figure~\ref{fig:tv7-task8}.
\end{itemize}

\begin{figure}[htbp]
  \centering
  \includegraphics[width=0.65\linewidth]{figures/TV7_task8.png}
  \caption{Task 8: \texttt{ImJPG} (left) and \texttt{ImJPG\_DeleteGreen}
    (right). The green channel becomes 0, so the green leaves turn dark purple.}
  \label{fig:tv7-task8}
\end{figure}

Figure~\ref{fig:tv7-singular-colors} shows two of these examples. After the
filter, the two colors in each row look exactly the same.

\begin{figure}[htbp]
  \centering
  \includegraphics[width=0.85\linewidth]{figures/TV7_singular_colors.png}
  \caption{Singular filters send different colors to the same color. Top: red and
    green both become the same gray with \texttt{GrayMatrix}. Bottom: two colors
    that differ only in green become the same color with
    \texttt{DeleteGreenMatrix}.}
  \label{fig:tv7-singular-colors}
\end{figure}

For these matrices \texttt{np.linalg.inv} gives an error, and no method can bring
back the lost channel. Some matrices are invertible but very close to singular,
like \texttt{SepiaMatrix}; they are studied in
Section~\ref{sec:theory-cond}.
